\documentclass[11pt]{article}
\usepackage[margin=2.6cm]{geometry}
\usepackage{graphicx,booktabs,amsmath,microtype,float}
\usepackage[hidelinks]{hyperref}
\newcommand{\code}[1]{\texttt{\small #1}}
\setlength{\parskip}{4pt}\setlength{\parindent}{0pt}
\title{A healthcare cost proxy from UKHLS alone\\[2pt]
\large What is buildable now, how to validate it, and when linkage becomes worth it}
\author{Julian Ashwin}
\date{September 2026}

\begin{document}
\maketitle

\noindent Companion to \code{health\_measures\_note.pdf} \S8. That section
established that expected cost, not a binary event, is the right outcome for
the framework's convexity question, and that UKHLS carries the physical
quantities cost is bought with. This is the plan for turning those quantities
into a defensible cost proxy without linkage.

\section{What we are building}

A per-person-wave estimate of annual NHS secondary-care cost, in a fixed
price year, for waves 7--15. It will be wrong in level and, we hope, roughly
right in \emph{shape} --- which is all the framework needs, since every
question we are asking is about how cost varies with health, not about its
absolute size.

The one-line target:
\[
  \widehat{C}_{it} \;=\;
  \underbrace{c_{\text{spell}}(k) + c_{\text{bed}} \cdot (\text{nights}_{it} - \bar{L})^{+}}_{\text{in-patient}}
  \;+\;
  \underbrace{c_{\text{op}} \cdot \widehat{n}^{\text{op}}_{it}}_{\text{out-patient}}
\]
with $k$ the attributed condition where we have one, and everything indexed
to a single price year.

\section{Ingredients we already hold}

\begin{center}\small
\begin{tabular}{llp{6.4cm}}
\toprule
Quantity & Source & Status \\
\midrule
Any admission & \code{hosp}, waves 7--15 & complete for respondents \\
Nights & \code{hospd} & all admitted; median 3, mean 7.5, max 300 \\
Maternity flag & \code{hospch} & 12.7\% of admissions, mean age 31.7 \\
Nights by new condition & \code{hospdc}$i$, waves 2--15 & 21\% of admissions \\
Nights by prior condition & \code{hospdcp}$i$, waves 11--15 & 11\% of admissions \\
\quad either & & 29\% of admissions, \textbf{65\% of nights} \\
Out-patient band & \code{hl2gp}, \code{hl2hop} & banded 0/1--2/3--5/6--10/$>$10 \\
\bottomrule
\end{tabular}
\end{center}

Unit costs come from the NHS National Cost Collection, which publishes
national average unit costs annually and free, at both aggregate and
HRG level. Indicative recent figures: non-elective in-patient episode
$\approx\pounds1{,}590$, elective $\approx\pounds3{,}684$, day case
$\approx\pounds738$, excess bed day $\approx\pounds313$, out-patient
attendance $\approx\pounds120$.

\section{Four stages}

\subsection*{Stage 1 --- a flat-rate baseline (half a day)}

Cost each admission at the non-elective episode rate plus excess bed days
beyond the average inclusive length of stay, and each out-patient attendance
at the flat rate, with band midpoints $0, 1.5, 4, 8, 15$. Exclude maternity.
This is deliberately crude and exists to give every later stage something to
be compared against.

\textbf{The one real modelling decision here.} UKHLS gives total nights but
\emph{not the number of spells}, so a person with 20 nights might have had one
long admission or four short ones, which differ in cost by a factor of about
three. Stage 1 assumes one spell. Stage 4 tests how much that matters; if the
answer is ``a lot'', it becomes the main argument for linkage, because HES
would resolve it exactly.

\subsection*{Stage 2 --- condition weighting (one to two days)}

For the $65\%$ of nights carrying an attribution, replace the flat rate with
a condition-specific one. The mapping is UKHLS condition code $\rightarrow$
NCC HRG chapter $\rightarrow$ average cost, at chapter level rather than
individual HRG, since the survey identifies a condition and not a procedure.
The attributed nights concentrate in stroke, COPD, asthma, osteoarthritis,
MI, congestive heart failure and chronic kidney disease, which span a wide
cost range, so this is where the proxy should improve most.

Unattributed nights take the case-mix-weighted average of the attributed
ones \emph{within the same health decile}, rather than a global average ---
otherwise the imputation flattens exactly the gradient we are trying to
measure.

\subsection*{Stage 3 --- validation against published profiles (one day)}

This is the stage that decides whether the proxy is usable, and it is the
reason to prefer building it here rather than assuming a shape.

\begin{enumerate}\itemsep2pt
\item \textbf{The age profile.} The CHE reference tables give average annual
NHS spend by age, sex and deprivation, and the IFS has published hospital
spending profiles from NHS administrative records. Our proxy should reproduce
the shape --- flat to about 50, rising steeply after 70 --- without having
been fitted to it. A proxy that misses the age profile cannot be trusted on
the health profile.
\item \textbf{The concentration.} Published figures put the top decile of
spenders at more than half of expenditure. Ours should land near that. Too
little concentration means the intensive margin is being under-weighted; too
much means outliers are running the estimate.
\item \textbf{The level.} Per-capita secondary care spend implied by the
proxy against national totals. We expect to \emph{undershoot}, because
UKHLS misses A\&E, and self-reported nights are known to be under-recalled;
the size of the shortfall is itself informative.
\end{enumerate}

\subsection*{Stage 4 --- use it, and test what it rests on (one to two days)}

Redo \S8's convexity analysis with $\pounds$ on the vertical axis, and
\S11's outcome prediction with cost as the outcome. Then the sensitivity
analysis that matters:

\begin{itemize}\itemsep2pt
\item \textbf{Spells.} Re-run assuming one spell, two spells, and
nights$/\bar{L}$ spells. If the convexity conclusion is stable across these,
the spell ambiguity is not binding and linkage buys less than it looks.
\item \textbf{Outlier handling.} The 300-night maximum will dominate any
mean. Winsorise at the 99th percentile and re-run.
\item \textbf{Condition weights.} Flat rate against condition-weighted: does
the shape change, or only the level?
\item \textbf{Maternity.} Already known to move the gradient from
$31.9\times$ to $52.7\times$; carry both through.
\end{itemize}

\section{What this cannot do}

Primary care, prescribing, A\&E and community services are absent entirely,
and they are not a small share of NHS spend. Out-patient use is a four-band
count, so its intensive margin is crude. Elective and emergency admissions
are indistinguishable, though they differ in cost by more than a factor of
two. And every quantity is self-reported over a twelve-month recall window,
against wave spacing that is not exactly twelve months.

The proxy is therefore a \emph{secondary-care bed-day-weighted cost index},
and should be named that way in anything we write rather than called a cost.

\section{The decision this sets up}

If Stage 3 shows the proxy reproducing the published age profile and
concentration, and Stage 4 shows the convexity conclusion stable across the
spell and outlier assumptions, then UKHLS alone supports the framework's
empirical claim and linkage is a refinement rather than a precondition.

If instead the shape moves with the spell assumption, or the proxy misses
the published profiles, then the case for UK LLC linkage becomes concrete
and specific --- and we would know exactly which quantity we needed it for,
which is a much stronger application than asking for HES in general. The
binding risk there remains sample size: roughly $8{,}000$ consented
participants against the $84{,}679$ in our current sample, with an extension
to all participants stated but needing verification.

\section*{Effort}

Three to five days end to end, entirely on data already on disk, with no new
permissions. Stages 1 and 3 are the ones worth doing even if we stop there:
a validated flat-rate index would already replace the physical-units analysis
of \S8 with something denominated in pounds.

\end{document}
